\documentclass[openany,fontset=fandol]{ctexbook}
\usepackage{amsmath,amssymb}
\usepackage[unicode,hidelinks]{hyperref}
\newcommand{\note}[1]{\footnote{译注：#1}}
\begin{document}
\title{代数几何基础第五卷预备笔记\\定理13.2及证明的初步归约}
\author{据用户上传英译编辑稿译出}
\date{2026年11月17日}
\maketitle
\chapter*{范围与校订状态}
本批连续译出底本第49页定理13.2的陈述及其后三段证明，自“Theorem 13.2 (M. Artin)”至“geometrically normal over $k$”。下一接点是第50页“The benefit of this hypothesis”。证明尚未译完，不包含引理13.3。

底本的术语不一致、引用疑问及归约中的略写均在脚注保留，不把本批当作独立完成的证明。未并入或覆盖主稿。
\chapter*{§13 初等态射与M. Artin定理（续）}
\textbf{定理13.2（M. Artin）。}设$k$是无限完美域，$X$是$k$上几何不可约的预概形，$x$是$X$的一个光滑点。则$x$具有一个由多初等开邻域组成的基本邻域系。\note{“多初等”沿用定义13.1，指结构态射为初等态射的复合。底本定理本句未逐项重述有限型等环境假设，而下文立即使用仿射、拟射影与射影闭包的归约；本批不将这些略写解读为对任意无有限性条件的预概形都已给出证明。}

把$X$替换为$x$的一个给定邻域后，只须证明$x$在$X$中有一个多初等开邻域〔底本英译者在句末注：最后这个字母是否正确？字迹不清〕。\note{底本此句写“open elementary neighborhood”，而定理陈述是“open polyelementary neighborhoods”，下一段又按维数归纳并构造$U\to V$，故此处正文统一为“多初等开邻域”，原用词留于此。初等结构态射的纤维为一维，不能直接把一般高维邻域称作在$k$上初等的。句末的字母难辨标记并未因此消除。}

对$X$的维数$n$作归纳，我们归约为证明：若$n>0$，则存在$x$的一个开邻域$U$及一个初等态射$f\colon U\to V$，其中$V$是$k$上的光滑概形（必然几何不可约，且维数为$n-1$）。（$n=0$的情形显然平凡，因为此时$X$同构于$\operatorname{Spec}(k)$；考虑到13.1并未排除空族态射的复合，它在$\operatorname{Spec}(k)$上是多初等的。）我们应当在13.1中以某种方式提到这一点。\note{此处保留原作者关于空复合的编写提醒。归纳中还须在$V$上取像点的适当邻域并拉回；底本没有在本段展开，本译文不把省略的归纳拼接冒充为原文已写出的步骤。}

\newpage
首先要求$k$完美的必要性，在$n=1$时就已显现：此时取$X$的函数域$K$的典范射影正规模型作为$X'$（参见第II章第7段）。\footnote{原脚注64：EGA II，是的，7.4〔英译者注〕。}$k$完美保证$X'$在$k$上光滑（因为$X'$无论如何都是正则的），并且也保证带诱导约化结构的$Z=X'\setminus X$在$k$上平展。现在来处理一般情形，于是可以假设$n\geq2$。\note{这一段关于曲线模型的写法隐含先在$x$附近作适当开限制，使所讨论的光滑曲线开集置于其正规射影模型中。并且定义13.1要求边界到基底满射，在域上即边界非空；若边界为空，还需在不含$x$处作适当删点。底本这里没有写出这些步骤，故指出略写而不声称已在正文补齐。}

显然可以假设$X$仿射，因而拟射影。再把$X$替换为一个射影闭包，就可以假设$X$射影，但始终须证明：每个邻域都包含一个允许初等态射$U\to V$的开邻域$U$。此外，把$X$替换为其正规化（它在$X$上有限，因而射影；引用待补），\footnote{原脚注65：我想是在EGA II〔英译者注〕。}并不会改变$x$的一个邻域，所以可以假设$X$正规；由于$k$完美，它因而在$k$上几何正规。\note{“不改变一个邻域”按局部含义译出，并非断言正规化在整个$X$上是同构。原稿对正规化有限性的引用尚待定位，原脚注65的疑问语气保留。本批只收入这些归约，不表示后续构造已经完成。}
\end{document}
